\section{Why an amplitude bound alone cannot settle the sparse case}
\label{sec:blocked-obstruction}
The remaining difficulty cannot be removed simply by proving that every
accurate low-degree threshold approximation to parity has amplitude growing
with the number of active coordinates. The following construction rules
out that assertion without any symmetry assumption on the polynomial.

\begin{proposition}[Bounded-amplitude blocked approximants]
\label{prop:blocked-approximants}
For every $\alpha,\varepsilon>0$, there are finite constants
$B=B(\alpha,\varepsilon)$ and $k_0=k_0(\alpha,\varepsilon)$ such that,
for every $d\ge k\ge k_0$, a multilinear polynomial $P$ on the $d$-cube
satisfies
\[
 \deg P\le\alpha k,\qquad
 1\le |P(x)|\le B\quad\hbox{for every }x,\qquad
 \Pr\{\operatorname{sign}P(X)\ne\chi_U(X)\}\le\varepsilon
\]
for any fixed $U\subseteq[d]$ of size $k$.
\end{proposition}
\begin{proof}
Replacing $\alpha,\varepsilon$ by smaller positive constants if necessary,
assume $\alpha\le1$ and $\varepsilon\le1/2$.
Choose a fixed integer $m$ sufficiently large that a central integer
interval $[A,B_0]$ for $T\sim\operatorname{Bin}(m,1/2)$ has probability
at least $1-\varepsilon$ and length $s=B_0-A\le\alpha m/2$.
Hoeffding's inequality permits $m=O(\alpha^{-2}\log(2/\varepsilon))$,
after increasing the absolute constant. Define
\[
 Q(t)=\sigma\prod_{j=A}^{B_0-1}(2t-2j-1),
\]
where $\sigma\in\{-1,1\}$ makes its sign agree with $(-1)^t$ throughout
$[A,B_0]\cap\mathbb Z$. At every integer $t\in[0,m]$, all factors are
nonzero odd integers, so $1\le|Q(t)|\le(2m+1)^s$.
Partition $U$ into $m$ nonempty blocks of size at most $\lceil k/m\rceil$
and let $Y_j$ be the parity of the $j$th block. These are independent
uniform signs, and their product is $\chi_U$. Set
\[
 P(x)=Q\!\left(\frac{m-\sum_{j=1}^mY_j(x)}2\right).
\]
Multilinearization in the block variables shows
$\deg P\le s\lceil k/m\rceil\le\alpha k$ whenever $k\ge m$.
The remaining conclusions follow from the bounds on $Q$ and the
binomial law of its argument.
\end{proof}

For example, if $3$ divides $k$, three equal blocks give
$P=-(Y_1+Y_2+Y_3)$, with degree $k/3$, global margin one, parity error
$1/4$, and $\|P\|_2^2=3$. If $4$ divides $k$, the polynomial
$P=1-\sum_{j=1}^4Y_j$ has degree $k/4$, error $5/16$, and
$\|P\|_2^2=5$. These examples need not be symmetric in all $k$ original
coordinates.

Proposition~\ref{prop:blocked-approximants} is an obstruction to a proof
method, not a learning upper bound for the posterior. If $\deg P<k$,
then $\mathbb E[\chi_U P]=0$; together with a global unit margin and
sign error $R$, this implies
\[
 \mathbb E(1-\chi_U P)_+^2\ge\frac1R.
\]
Thus an accurate blocked approximant has larger squared repair loss
than the constant polynomial, whose loss is exactly two. Determining
the minimum energy of repairs that also preserve the global margin
remains essential.
